% Generated draft: requires mathematical and source review.
\documentclass[11pt]{article}
\usepackage[margin=25mm]{geometry}
\usepackage{fontspec}
\setmainfont{Latin Modern Roman}
\usepackage{amsmath,amssymb,mathtools,unicode-math}
\setmathfont{Latin Modern Math}
\usepackage{enumitem,needspace,xurl,hyperref}
\setlength{\parindent}{0pt}
\setlength{\parskip}{5pt}
\setlength{\emergencystretch}{3em}
\newcommand{\N}{\mathbb{N}}
\newcommand{\Z}{\mathbb{Z}}
\newcommand{\Q}{\mathbb{Q}}
\newcommand{\R}{\mathbb{R}}
\newcommand{\C}{\mathbb{C}}
\newcommand{\F}{\mathbb{F}}
\newcommand{\A}{\mathbb{A}}
\newcommand{\PP}{\mathbb{P}}
\newcommand{\E}{\mathbb{E}}
\newcommand{\eps}{\varepsilon}
\newcommand{\dd}{\,\mathrm{d}}
\begin{document}
\section*{UnsolvedMath ID: 6700024}
\textbf{Source problem code:} catalogue entry
\subsection{Definitions and mathematical statement}
A group $\Gamma$ is residually finite if
\[
 \forall\gamma\in\Gamma\setminus\{1\}\quad
 \exists Q\text{ finite}\quad\exists\phi:\Gamma\to Q
 \text{ a homomorphism with }\phi(\gamma)\ne1.
\]
Equivalently the intersection of its finite-index normal subgroups is trivial. For a manifold $X$, the source asks whether residual finiteness of its fundamental group $\pi_1(X)$ is essential.

The theorem or construction to which ``essential'' refers is absent from the supplied record. Consequently neither the remaining manifold hypotheses nor the conclusion whose dependence is being tested is specified. To ask whether a residual-finiteness hypothesis can be removed, or whether a counterexample exists without it, requires that antecedent implication first be identified.

The research summary suggests scalar-curvature, index/Novikov, aspherical, covering, or assembly arguments as possible context, but does not supply an exact statement. No one of those conjectures or theorems is selected as the target. The retained problem is precisely the source's hypothesis-necessity question, with the missing antecedent recorded as a blocker to a fully quantified formulation.
\subsection{Short English statement}
Is residual finiteness of the fundamental group necessary for the result under discussion in the original source? The relevant result is missing, so no exact removal-of-hypothesis claim can be stated. The supplied summary reports no direct resolution but gives only possible scalar-curvature/index context. Its OPEN-TRIAGE label and the historical 2017 question are retained with current status needing review.
\subsection{Sources}
\begin{itemize}
\item Local UnsolvedMath registry, permanent ID 6700024, source code catalogue entry, scalar-curvature Question [?23], proposed by Misha Gromov, 2017. Dataset curation and metadata: Ulam AI and UnsolvedMath Contributors, CC BY 4.0.
\item Misha Gromov, \emph{101 Questions, Problems and Conjectures around Scalar Curvature} (2017), printed tag [?23], occurrence 24, PDF page 34, \url{https://www.ihes.fr/~gromov/wp-content/uploads/2018/08/101-problemsOct1-2017.pdf}; normalized pdf-text extraction, accessed 2026-07-29.
\item The supplied record contains neither the antecedent result nor an exact later primary citation. Source-specific provenance: the author-hosted PDF is incomplete and publicly available; current status and statement-reuse terms NEEDS\_REVIEW, separately from the dataset licence.
\end{itemize}
\end{document}
